\section{Introduction}\label{sec:introduction}

Grothendieck's homotopy hypothesis asks whether weak globular
infinity-groupoids describe the homotopy theory of spaces.  In
\emph{Pursuing Stacks}, the operations on paths and higher homotopies
are specified by freely adjoining coherence operations to globular
pasting diagrams \cite[Sections~9 and~12--13]{Grothendieck1983}.
A \emph{coherator} records these choices.  Maltsiniotis gave a precise
formulation in terms of coherators and their models, together with
realization and fundamental-infinity-groupoid functors
\cite{Maltsiniotis2007,Maltsiniotis2010}.
The resulting categories depend on the chosen coherator; the hypothesis
asserts that their homotopy theory recovers that of spaces.
We establish it for every coherator in the Ara--Henry convention
specified below.

The choice of algebraic model is substantive.  Simplicial sets already
model spaces through the classical realization--singular-complex
Quillen equivalence \cite[Theorem~3.6.7]{Hovey1999}.
For strict globular infinity-groupoids, however, the simply connected
homotopy types represented under a realization preserving homotopy
groups are precisely products of Eilenberg--Mac~Lane spaces \cite[Theorem~5.4]{AraStrict2013}.
The passage from strict identities to weak coherence is therefore
substantive.  Here weak equivalences are defined
inside the chosen globular algebraic category, using components and
based homotopy groups.

Henry isolated a concrete obstruction to constructing their canonical
semi-model structure \cite[Conjecture~5.3.3]{Henry2016}.  Suppose that \(X\)
is built by attaching globular cells.  Select an \(n\)-cell \(a\) of \(X\),
adjoin a new parallel \(n\)-cell \(a'\), and adjoin an \((n+1)\)-cell from
\(a\) to \(a'\).  The conjecture asserts that this free extension is a weak
equivalence.  Although the extension has an elementary retraction, its
free algebraic construction also creates composites and higher coherence
cells.  Controlling the homotopy classes of all those cells is the issue.

We use Grothendieck coherators in the convention of Ara
\cite[Sections~2.3 and~2.9--2.11]{Ara2013}, fixed by Henry in
\cite[Section~5.3]{Henry2016}; the precise definition appears in
Section~\ref{sec:preliminaries}.  Thus generating fillers of output
dimension \(j\) have pasting-scheme arity of dimension at most \(j\).
Write \(D_j\) for the free model on a \(j\)-globe and \(\partial D_j\)
for its boundary, with \(\partial D_0=\varnothing\).
A model is \emph{cellular} if it is obtained from the initial model by
a possibly transfinite sequence of boundary attachments.
Weak equivalences induce a bijection on components and isomorphisms on
all based homotopy groups.

\begin{samepage}
\begin{theorem}\label{thm:main}
Fix any Grothendieck coherator \(\mathcal C\).
Let \(X\) be a cellular \(\mathcal C\)-infinity-groupoid, let \(n\ge0\),
and let \(a:D_n\to X\) be any cell.  In the pushout square
\[
\begin{tikzpicture}[baseline=(current bounding box.center),
                   node distance=1.25cm and 2.2cm,>=Stealth]
\node (a) {\(D_n\)};
\node (b) [right=of a] {\(D_{n+1}\)};
\node (c) [below=of a] {\(X\)};
\node (d) at (b |- c) {\(X^+\)};
\draw[->] (a) -- node[above] {\(s_n\)} (b);
\draw[->] (a) -- node[left] {\(a\)} (c);
\draw[->] (b) -- (d);
\draw[->] (c) -- node[below] {\(i\)} (d);
\end{tikzpicture}
\]
where \(s_n\) is the source inclusion, the map \(i:X\to X^+\)
is a weak equivalence.
\end{theorem}
\end{samepage}

The pushout is taken in models of the fixed theory \(\mathcal C\).
The selected cell \(a\) can be any composite or coherence cell.
For \(n=0\), the extension adjoins an object and an arrow from the selected
object to it.  In particular, Theorem~\ref{thm:main} resolves Henry's
pushout conjecture affirmatively.  Henry already showed how its
finitely cellular case extends to transfinite cellular objects and their
cofibrant retracts \cite[Lemma~5.3.4]{Henry2016}; the proof here treats
arbitrary cellular objects directly.
The canonical semi-model structure and its equivalence with the homotopy
theory of spaces follow in Section~\ref{sec:consequences}.

The algebraic and homotopical sides of this question developed together.
Berger's globular theories organize operations by tree-shaped pasting
arities \cite[Definition~1.5]{Berger2002}.
Ara constructed the homotopy-group calculus and fixed-boundary tests
for weak equivalences used here \cite[Section~4]{Ara2013}.
Bourke described cellular globular theories through iterated algebraic
injectivity and proved faithfulness of their free stages
\cite[Theorems~3.12 and~5.2]{Bourke2020}.
These results make the algebraic models and their weak equivalences
precise; the remaining task is to show that elementary expansions have
the required homotopical behavior.

Henry identified the elementary pushout assertion as sufficient both
for the canonical semi-model structure and for the comparison with
spaces \cite[Theorem~5.3.5 and Corollary~5.3.13]{Henry2016}.
One route to this assertion is a functorial path construction.
Bourke's endomorphism-theory argument puts weak groupoid structures
on iterated path objects in identity type categories
\cite[Theorem~5.1]{Bourke2016}.
Within Grothendieck groupoids, Lanari constructed a globular candidate
with compositions, identities and inverses
\cite[Definitions~5.2 and~6.1, Proposition~6.2, and Theorem~6.5]{Lanari2020}.
His construction leaves the extension to compatible interpretations
of all higher operations open.  The partial paths used here are
adapted to one expansion index at a time: the already established
higher-index expansions make their representing objects contractible,
which then permits every operation of the fixed coherator to be
interpreted in its original free order.

In dimension three, Lanari obtained the canonical semi-model structure
for Grothendieck groupoids of a suitable type
\cite[Theorem~4.10]{Lanari2018}, and Henry and Lanari proved the
corresponding homotopy-hypothesis comparison \cite{HenryLanari2023}.
Their study of coskeletal and truncated globular models also supplies
context for the cutoffs used below \cite[Section~2]{HenryLanari2023}.
Taylor's later construction of algebraic coherators and Grothendieck
realizations retains a generalized pushout conjecture as a hypothesis
\cite[Conjecture~5.7 and Theorem~5.8]{Taylor2026}.
Theorem~\ref{thm:main} proves Henry's assertion for an arbitrary fixed
coherator in our convention.  Its consequence in
Section~\ref{sec:consequences} is a Quillen equivalence with spaces,
and hence independence of the resulting homotopy theory from that
choice of coherator.

\subsection*{The proof}

Section~\ref{sec:preliminaries} first expresses weak equivalence as a
lifting test with a prescribed boundary.  The proof can then work with
individual cells and their joining homotopies throughout.
The next step, in Section~\ref{sec:cutoffs-trees}, is a reflection to models in which every boundary above
dimension \(q\) has a unique filler.  It preserves all cells through \(q\).
We then hold \(q\) fixed and descend in the index \(n\) of the elementary
expansion.  The top case is immediate from the retraction and the unique
higher fillers.

At the induction step, expansions in dimensions greater than \(n\)
are already known to be weak equivalences.  A factorization argument
then gives Lemma~\ref{lem:upper-cell}: every weak equivalence built by boundary attachments
strictly above \(n\) splits and remains a weak equivalence after the
cellular pushouts needed in the proof.

The main geometric input concerns finite trees in a globular
pasting diagram.  Such a tree can be compared with a string on the same
vertices.  The composite comparison is homotopic to the identity on its
generating edges.  To pass from these individual homotopies
to a statement about all coherences, we duplicate the top generators
using only expansions already covered by the induction.
Freeness and two-out-of-three then prove contractibility of the whole
tree model (Proposition~\ref{prop:tree-contractible}).

In Section~\ref{sec:partial-cylinders}, these tree models provide partial
cylinders: an \(n\)-cell moves along an \((n+1)\)-cell, while its
lower boundary stays fixed.  Higher cylinders fit these movements
together along shared faces.  A skeletal gluing formula proves that
the objects representing such compatible tuples of cylinders are
contractible.  Interpreting the freely adjoined operations in their
chosen stage order then makes the associated path construction a functor on
models of \(\mathcal C\).  Its two endpoint maps record the ends of
each cylinder.  Individually they lift all boundary inclusions; together
they lift the current source inclusion (Proposition~\ref{prop:partial-path}).
The endpoint maps produce self-maps of the source and target that are
weak equivalences; a retract argument then proves that the expansion
itself is a weak equivalence.  Each fixed-boundary
test uses finitely many dimensions, so the end of
Section~\ref{sec:partial-cylinders} recovers
both the tested disk and its homotopy witness before the cutoff.
Section~\ref{sec:consequences} derives the canonical left semi-model
structure and applies Henry's comparison with spaces.

The main construction uses three orders: at fixed \(q\), descending
induction on the expansion index \(n\); within each step, increasing
cylinder dimension \(j\); and, for operation interpretations, the
chosen free stage order of \(\mathcal C\).
Its reusable ingredients are the reflection
preserving lower cells, the promotion of generator homotopies by
higher-dimensional expansions, and the use of partial cylinders tailored
to a single expansion index.
